% GENERATED FILE. Edit canonical lessons, then run npm run generate:books.
% canonical-source-hash: fnv1a64:bfb69cba2ce9af7f
% canonical-lessons: UR-C04-kaise-kaisi, UR-W04-joining, UR-C04-aap-kaise-hain, UR-C04-main-hun, UR-W04-kaf, UR-C04-thik, UR-C04-main-thik-hun, UR-W04-nun, UR-C04-practice

\chapter{How Are You?}
\label{ch:ur-how-are-you}
\hlchaptermodality{4}

\begin{chapteropening}
\textbf{By the end of this chapter:} \emph{I can ask how someone is with respectful āp, choose kaise or kaisī for the person I am addressing, and reply that I am fine, thank you. I can join four letters into salām, and write kāf and nūn.}

The wellbeing exchange run both ways — āp kaise haiṅ? and āp kaisī haiṅ? — answered with the unchanging maiṅ ṭhīk hūṅ, shukriyā.
\end{chapteropening}

\section[kaise / kaisī]{\ur{کیسے} / \ur{کیسی} — “how” follows the addressee}

\label{lesson:UR-C04-kaise-kaisi}



\begin{quote}
Say \textbf{kyā} “what.” Urdu's next question word begins with the same \textbf{k} family, but its ending changes with the person being addressed.
\end{quote}

\subsection*{The two words}
\noindent
\begin{tabularx}{\linewidth}{@{}>{\raggedright\arraybackslash}X>{\raggedright\arraybackslash}X>{\raggedright\arraybackslash}X@{}}
\toprule
\textbf{Urdu} & \textbf{read} & \textbf{use in the respectful question} \\
\midrule
\textbf{\ur{کیسے}} & \emph{kaise} & addressing a man \\
\textbf{\ur{کیسی}} & \emph{kaisī} & addressing a woman \\
\bottomrule
\end{tabularx}

Both begin \textbf{\ur{ک}} \emph{k} + \textbf{\ur{ی}} carrying the \emph{ai} sound. At the left edge, \textbf{\ur{کیسے}} ends in broad \textbf{\ur{ے}} \emph{e}, while \textbf{\ur{کیسی}} ends in \textbf{\ur{ی}} long \emph{ī}. Read the whole word before isolating the final shape.

\begin{grammarlens}[title={Grammar Lens: one agreement contrast}]
English keeps \emph{how} unchanged. Urdu makes this describing form agree with the person addressed: \textbf{kaise} for a man, \textbf{kaisī} for a woman. Learn only this pair today—not a table of adjective endings. Respectful \textbf{āp} itself remains the same for everyone.
\end{grammarlens}

\subsection*{Your turn}
\begin{itemize}
\raggedright
  \item \emph{Say it:} addressing a man — \textbf{kaise}
  \item \emph{Say it:} addressing a woman — \textbf{kaisī}
  \item \emph{Point to:} broad final \textbf{\ur{ے}} versus long-\emph{ī} \textbf{\ur{ی}}
  \item \emph{Recall:} say \emph{nahī̃}
\end{itemize}

\begin{tcolorbox}[breakable,colback=teal!4,colframe=teal!35!black,title={Before you move on}]
Which form addresses a man? (\textbf{Kaise}.) A woman? (\textbf{Kaisī}.)

Source: \href{https://openbooks.lib.msu.edu/urdu/chapter/2-4/}{\emph{Basic Urdu}: Formal Conversation}.
\end{tcolorbox}
\section[salām]{Letters that hold hands, and the one that will not}

\label{lesson:UR-W04-joining}



\begin{quote}
Write \ur{س}, then \ur{م}. Four letters are now yours.
\end{quote}

\subsection*{Script — the running, falling line}
Urdu letters \textbf{join}. They run into one another along a single connected line, and because this is Nastaliq that line is also \textbf{falling} — each joined letter sits a little lower and further left than the one before it.

Because they join, most letters \textbf{change shape} by position. A letter opening a joined run, one buried in the middle, one closing it, and one standing alone can all look different. They are the same letter.

That sounds like four times the work and it is not, because the change is almost always \textbf{the same trimming}: the part that identifies the letter stays, and the finishing flourish is cut short so the next letter can start. You saw it already in mīm, where joining costs the tail and keeps the head.

\subsection*{Script — and the letter that refuses}
Now the rule that saves the most confusion.

\textbf{Some letters take the hand of the letter on their right and offer nothing to the left.} The run stops dead after them, and the next letter begins as though starting a fresh word.

\textbf{alif is one of these.} So a gap in the middle of an Urdu word is usually not a word break at all — it is a non-joining letter doing exactly what it always does. Without that fact a page of Urdu looks like it has been spaced at random.

\begin{grammarlens}[title={What you've built}]
Four letters, and they spell the word this book opened with:

\begin{quote}
\ur{سلام}
\end{quote}

From the right: \textbf{\ur{س}} \emph{s}, \textbf{\ur{ل}} \emph{l}, \textbf{\ur{ا}} \emph{ā}, \textbf{\ur{م}} \emph{m}. Say it: \textbf{salām}.

Two things to see. The short vowel between s and l is \textbf{not written} — Urdu prints the long vowels and leaves the short ones to the reader. And the break before the final \ur{م} is alif refusing to join, exactly as promised.
\end{grammarlens}

\subsection*{Your turn}
\begin{itemize}
\raggedright
  \item \emph{Read it:} \ur{سلام} — name each letter from the right, then say the word
  \item \emph{Write it:} \ur{سلام} — let the line fall as it goes, and let it break after alif
  \item \emph{Say it:} \textbf{salām}
\end{itemize}

\begin{tcolorbox}[breakable,colback=teal!4,colframe=teal!35!black,title={Before you move on}]
What does mīm lose when it joins? (Its \textbf{tail}.)
\end{tcolorbox}
\section[āp kaise haiṅ? / āp kaisī haiṅ?]{\ur{آپ} \ur{کیسے} \ur{ہیں؟} / \ur{آپ} \ur{کیسی} \ur{ہیں؟} — ask respectfully}

\label{lesson:UR-C04-aap-kaise-hain}



\begin{quote}
Choose respectful \textbf{āp}. Then choose \textbf{kaise} for a man or \textbf{kaisī} for a woman. One final copula completes the question.
\end{quote}

\subsection*{The exchange}
\begin{quote}
to a man: \textbf{\ur{آپ} \ur{کیسے} \ur{ہیں؟}} — \emph{āp kaise haiṅ?} to a woman: \textbf{\ur{آپ} \ur{کیسی} \ur{ہیں؟}} — \emph{āp kaisī haiṅ?}
\end{quote}

The order is \textbf{you + how + are?} Keep \textbf{\ur{ہیں}} \emph{haiṅ} final, just as \textbf{\ur{ہے}} \emph{hai} stayed final in the name question. Respectful \textbf{āp} takes this plural-shaped copula even when speaking to one person.

\begin{grammarlens}[title={Grammar Lens: respect does not erase agreement}]
Both lines are respectful. The choice between \textbf{kaise} and \textbf{kaisī} tracks the person addressed, not greater or lesser politeness. If you do not know how a person wishes to be addressed, listen for their usage or ask rather than guess; the grammar lesson does not replace social judgment.
\end{grammarlens}

\subsection*{Your turn}
\begin{itemize}
\raggedright
  \item \emph{Say it:} to a man — \textbf{āp kaise haiṅ?}
  \item \emph{Say it:} to a woman — \textbf{āp kaisī haiṅ?}
  \item \emph{Keep:} respectful \textbf{āp} and final \textbf{haiṅ} in both
\end{itemize}

\begin{tcolorbox}[breakable,colback=teal!4,colframe=teal!35!black,title={Before you move on}]
What stays the same in both lines? (\textbf{Āp ... haiṅ}.) What changes? (\textbf{Kaise / kaisī}.)

Sources: \href{https://openbooks.lib.msu.edu/urdu/chapter/2-4/}{\emph{Basic Urdu}: Formal Conversation} and \href{https://openbooks.lib.msu.edu/urdu/chapter/2-7/}{Personal Pronouns with “To Be”}.
\end{tcolorbox}
\section[maiṅ ... hūṅ]{\ur{میں} ... \ur{ہوں} — one frame for “I am ...”}

\label{lesson:UR-C04-main-hun}



\begin{quote}
Ask \textbf{āp kaise haiṅ?} You know the respectful “you are” ending \textbf{haiṅ}. Your answer needs the first-person pair “I ... am.”
\end{quote}

\subsection*{You'll want to know first — the frame}
\begin{quote}
\textbf{\ur{میں} ... \ur{ہوں}} — \emph{maiṅ ... hūṅ} — \textbf{I am ...}
\end{quote}

\textbf{\ur{میں}} \emph{maiṅ} means “I.” Final \textbf{\ur{ں}} marks a nasal ending. \textbf{\ur{ہوں}} \emph{hūṅ} means “am” and carries long \emph{ū} with \textbf{\ur{و}}, followed again by nasal \textbf{\ur{ں}}. Leave a space for the state word that comes next.

\begin{grammarlens}[title={Grammar Lens: the verb closes the frame}]
Urdu puts the copula at the end:

\begin{quote}
\textbf{maiṅ + {[}state{]} + hūṅ} — I + {[}state{]} + am
\end{quote}

Learn only this first-person frame. You do not need the full copula chart to answer one wellbeing question. Notice the useful contrast: respectful \textbf{āp ... haiṅ}, but first-person \textbf{maiṅ ... hūṅ}.
\end{grammarlens}

\subsection*{Your turn}
\begin{itemize}
\raggedright
  \item \emph{Say it:} \textbf{maiṅ} — I
  \item \emph{Say it:} \textbf{hūṅ} — am
  \item \emph{Frame:} \textbf{maiṅ ... hūṅ}, leaving the middle open
\end{itemize}

\begin{tcolorbox}[breakable,colback=teal!4,colframe=teal!35!black,title={Before you move on}]
Which word opens the answer? (\textbf{Maiṅ}.) Which copula closes it? (\textbf{Hūṅ}.)

Source: \href{https://openbooks.lib.msu.edu/urdu/chapter/2-7/}{\emph{Basic Urdu}: Personal Pronouns with “To Be”}.
\end{tcolorbox}
\section[kāf]{\ur{ک} — a bowl with a slash laid over it}

\label{lesson:UR-W04-kaf}



\begin{quote}
Write \ur{سلام}. Four letters, one break.
\end{quote}

\subsection*{Script — the shape}
\begin{quote}
\ur{ک}
\end{quote}

Two movements, and exactly \textbf{one} lift between them.

\begin{enumerate}
\raggedright
  \item a \textbf{stem} down, flowing right to left through a \textbf{flatter bowl} and finishing with a pronounced \textbf{hook} — all without lifting
  \item lift once, then a \textbf{long slash} drawn down from the upper right towards the stem
\end{enumerate}

The slash is not a dot and not decoration; it is part of the letter, and leaving it off writes nothing at all.

Its name is \textbf{kāf} and it says \textbf{k}.

Worth knowing: \textbf{this is the Urdu kāf.} Arabic writes the same sound with a different shape, and if you have seen Arabic you will notice the difference immediately. Urdu keeps the Perso-Arabic alphabet and does not keep every Arabic shape.

\subsection*{Your turn}
\begin{itemize}
\raggedright
  \item \emph{Write it:} \ur{ک} — stem, bowl and hook in one movement, then lift and lay the slash across
  \item \emph{Say it:} \textbf{k}
\end{itemize}

\begin{tcolorbox}[breakable,colback=teal!4,colframe=teal!35!black,title={Before you move on}]
Is the slash optional? (\textbf{No} — without it the letter is not written.)
\end{tcolorbox}
\section[ṭhīk]{\ur{ٹھیک} — “fine,” “well,” and “correct”}

\label{lesson:UR-C04-thik}



\begin{quote}
Open and close the answer frame: \textbf{maiṅ ... hūṅ}. The middle now gets one high-frequency state word.
\end{quote}

\subsection*{You'll want to know first — the word}
\begin{quote}
\textbf{\ur{ٹھیک}} — \emph{ṭhīk} — \textbf{fine, well, correct}
\end{quote}

The first sound is retroflex and aspirated: curl the tongue slightly back for \textbf{\ur{ٹ}} \emph{ṭ}, then release a puff marked by \textbf{\ur{ھ}}. A plain English \emph{t} is a useful first approximation while the contrast settles. \textbf{\ur{ی}} carries long \emph{ī} before final \textbf{\ur{ک}} \emph{k}.

\begin{cousinweb}
Urdu \textbf{\ur{ٹھیک}} and Hindi \emph{ṭhīk} are the same everyday Hindustani word. Their grammar is closely shared; their scripts make different literary histories visible. This Urdu lesson leaves the Hindi side romanized: mixed-language practice may show both scripts only after each form has been learned in its own right.
\end{cousinweb}

\subsection*{Your turn}
\begin{itemize}
\raggedright
  \item \emph{Say it:} \textbf{ṭhīk} — fine, well
  \item \emph{Feel:} tongue slightly back, then a small puff
  \item \emph{Connect:} Urdu \textbf{\ur{ٹھیک}} $\leftrightarrow$ Hindi \emph{ṭhīk}
  \item \emph{Recall:} say \emph{āp / tum / tū}, then say \emph{kyā}
\end{itemize}

\begin{tcolorbox}[breakable,colback=teal!4,colframe=teal!35!black,title={Before you move on}]
Where will \textbf{ṭhīk} go in the frame \textbf{maiṅ ... hūṅ}? (In the middle.)

Source: \href{https://openbooks.lib.msu.edu/urdu/chapter/2-2/}{\emph{Basic Urdu}: Greetings and Introductions}.
\end{tcolorbox}
\section[maiṅ ṭhīk hūṅ, shukriyā]{\ur{میں} \ur{ٹھیک} \ur{ہوں،} \ur{شکریہ} — “I am fine, thank you”}

\label{lesson:UR-C04-main-thik-hun}



\begin{quote}
Fill the frame \textbf{maiṅ ... hūṅ} with \textbf{ṭhīk}, then recall Chapter 1's \textbf{shukriyā}.
\end{quote}

\subsection*{The exchange}
\begin{quote}
\textbf{\ur{میں} \ur{ٹھیک} \ur{ہوں،} \ur{شکریہ۔}} \emph{Maiṅ ṭhīk hūṅ, shukriyā.} \textbf{I am fine, thank you.}
\end{quote}

No word is new. The step is assembling the words quickly enough to answer a real person.

\begin{grammarlens}[title={Grammar Lens: keep the copula final}]
The core reply is:

\begin{quote}
I + fine + am $\to$ \textbf{maiṅ ṭhīk hūṅ}
\end{quote}

English places “am” after “I”; Urdu closes the clause with \textbf{hūṅ}. \textbf{Shukriyā} comes after the completed answer as the courtesy move you already own.
\end{grammarlens}

\subsection*{Your turn}
\begin{itemize}
\raggedright
  \item \emph{Build:} \textbf{maiṅ + ṭhīk + hūṅ}
  \item \emph{Add:} \textbf{shukriyā}
  \item \emph{Say it:} the full reply without translating each word
\end{itemize}

\begin{tcolorbox}[breakable,colback=teal!4,colframe=teal!35!black,title={Before you move on}]
Which word must close the core answer? (\textbf{Hūṅ}.)

Source: \href{https://openbooks.lib.msu.edu/urdu/chapter/2-4/}{\emph{Basic Urdu}: Formal Conversation}.
\end{tcolorbox}
\section[nūn]{\ur{ن} — a deep bowl, and one dot}

\label{lesson:UR-W04-nun}



\begin{quote}
Write \ur{ک}. Bowl and hook, lift, slash.
\end{quote}

\subsection*{Script — the shape}
\begin{quote}
\ur{ن}
\end{quote}

A \textbf{bowl} swept from right to left, dropping well \textbf{below the baseline} and rising again — then one lift, and a single \textbf{dot} placed near the baseline, inside the curve.

Two parts, one lift. The bowl is deeper and rounder than the flatter one you drew in kāf, and the dot sits low rather than floating above.

Its name is \textbf{nūn} and it says \textbf{n}.

\subsection*{Your turn}
\begin{itemize}
\raggedright
  \item \emph{Write it:} \ur{ن} — the deep bowl, then lift and place the dot low inside it
  \item \emph{Say it:} \textbf{n}
\end{itemize}

\begin{tcolorbox}[breakable,colback=teal!4,colframe=teal!35!black,title={Before you move on}]
Which bowl is deeper, nūn's or kāf's? (\textbf{nūn's}.)
\end{tcolorbox}
\section[Practice]{Practice — an Urdu wellbeing exchange}

\label{lesson:UR-C04-practice}



\begin{quote}
Choose the question ending for the person addressed, then retrieve the complete answer before looking below.
\end{quote}

\subsection*{The exchange}
\begin{quote}
A to a man: \textbf{\ur{سلام}! \ur{آپ} \ur{کیسے} \ur{ہیں؟}} \emph{Salām! Āp kaise haiṅ?} A to a woman: \textbf{\ur{سلام}! \ur{آپ} \ur{کیسی} \ur{ہیں؟}} \emph{Salām! Āp kaisī haiṅ?} B: \textbf{\ur{میں} \ur{ٹھیک} \ur{ہوں،} \ur{شکریہ۔}} \emph{Maiṅ ṭhīk hūṅ, shukriyā.}
\end{quote}

The question keeps respectful \textbf{āp} and final \textbf{haiṅ} while \textbf{kaise/kaisī} agrees with the addressee. The reply is unchanged for the speaker's gender: \textbf{ṭhīk} does not change, and first-person \textbf{hūṅ} closes the line.

\begin{grammarlens}[title={Grammar Lens: separate grammar from guessing}]
The two question forms are a grammatical contrast, not an instruction to infer someone's identity. In real interaction, follow the person's own language or ask. The stable beginner default is respect: \textbf{āp ... haiṅ}.
\end{grammarlens}

\subsection*{Your turn}
\begin{itemize}
\raggedright
  \item \emph{Say it:} the man-addressed exchange
  \item \emph{Say it:} the woman-addressed exchange
  \item \emph{Answer:} \textbf{maiṅ ṭhīk hūṅ, shukriyā} after either question
\end{itemize}

\begin{tcolorbox}[breakable,colback=teal!4,colframe=teal!35!black,title={Before you move on}]
Run both question variants once. If the choice stalls, review only \textbf{kaise/kaisī}; if the answer order stalls, review only \textbf{maiṅ ... hūṅ}.
\end{tcolorbox}
